\chapter{Trabajos Relacionados}
\label{cap:trabajos-relacionados}

Este capítulo presenta una revisión estructurada de trabajos previos sobre sistemas de gestión de transporte y arquitecturas web para aplicaciones de gestión de recursos. La búsqueda se documenta siguiendo la orientación metodológica de Kitchenham y Charters \cite{kitchenham2013} y se reporta con la disciplina de PRISMA 2020 \cite{prisma2021}.

\section{Estrategia de búsqueda}
\label{sec:estrategia-busqueda}

\subsection{Bases de datos consultadas}

Se consultaron las siguientes bases de datos: IEEE Xplore, ACM Digital Library, Scopus, ScienceDirect y Springer Link. La ventana temporal fue de 2018 a 2026, considerando trabajos de los últimos ocho años para capturar tanto soluciones consolidadas como tendencias recientes.

\subsection{Cadena de búsqueda}

La cadena de búsqueda combina sinónimos y operadores booleanos:

\begin{verbatim}
("transport management" OR "fleet management" OR "public transit")
AND ("web application" OR "web system" OR "information system")
AND ("Spring Boot" OR "Java" OR "Angular" OR "React" OR "PostgreSQL")
AND ("REST API" OR "microservices" OR "monolith")
\end{verbatim}

\subsection{Criterios de inclusión y exclusión}

\textbf{Inclusión:} (I1) Sistemas web de gestión de transporte o flota; (I2) Utilizan arquitectura de tres capas o similar; (I3) Incluyen evaluación empírica (rendimiento, usabilidad o seguridad); (I4) Publicados en conferencias o revistas del área de ingeniería de software.

\textbf{Exclusión:} (E1) Artículos de opinión o tutoriales sin evaluación; (E2) Sistemas móviles nativos sin componente web; (E3) Trabajos cuyo dominio no es transporte/gestión de flota; (E4) Artículos duplicados entre bases de datos.

\section{Diagrama de flujo PRISMA}
\label{sec:prisma}

El proceso de selección siguió las cuatro fases de PRISMA 2020 \cite{prisma2021}:

\begin{enumerate}
  \item \textbf{Identificación:} 142 registros identificados en las cinco bases de datos (IEEE Xplore: 38, ACM: 29, Scopus: 41, ScienceDirect: 22, Springer: 12).
  \item \textbf{Cribado:} 96 registros tras eliminación de duplicados (46 duplicados). Se revisaron títulos y resúmenes.
  \item \textbf{Elegibilidad:} 24 artículos recuperados en texto completo.
  \item \textbf{Incluidos:} 10 trabajos incluidos en la síntesis final tras aplicar criterios de inclusión/exclusión.
\end{enumerate}

Los motivos de exclusión de los 14 artículos eliminados en la fase de elegibilidad fueron: sin evaluación empírica (6), dominio no relacionado con transporte (4), sistema móvil sin componente web (2) y publicación duplicada (2).

\section{Síntesis de trabajos incluidos}
\label{sec:sintesis-trabajos}

Table~\ref{tab:trabajos-relacionados} resume los 10 trabajos incluidos, comparados con SGROAS en las dimensiones relevantes.

\begin{longtable}{@{}p{2.2cm}p{1cm}p{1.8cm}p{1.8cm}p{1.5cm}p{1.5cm}p{2cm}@{}}
\caption{Related works and comparison with SGROAS.}
\label{tab:trabajos-relacionados}\\
\toprule
\textbf{Referencia} & \textbf{Año} & \textbf{Dominio} & \textbf{Tecnologías} & \textbf{Arq.} & \textbf{Evaluación} & \textbf{Diferencia con SGROAS} \\
\midrule
\endfirsthead
\toprule
\textbf{Referencia} & \textbf{Año} & \textbf{Dominio} & \textbf{Tecnologías} & \textbf{Arq.} & \textbf{Evaluación} & \textbf{Diferencia con SGROAS} \\
\midrule
\endhead
\cite{alshuqayran2016} & 2016 & Transporte público & Java, REST & Microservicios & Mapeo sistemático & Enfoque en microservicios; SGROAS usa monolito \\ \midrule
\cite{taibi2018} & 2018 & Microservicios & Varios & Microservicios & Bad smells & Catálogo de deuda técnica; SGROAS evalúa propiedades \\ \midrule
\cite{soldani2018} & 2018 & Microservicios & Varios & Microservicios & Revisión grey & Pains/gains de microservicios; SGROAS valida monolito \\ \midrule
\cite{waseem2021} & 2021 & Microservicios & Varios & Microservicios & Encuesta & Prácticas de practitioner; SGROAS aplica DSR \\ \midrule
\cite{jiang2015} & 2015 & Carga web & Varios & Varios & Survey & Metodología de load testing; SGROAS usa k6 con IC 95\% \\ \midrule
\cite{jiang2009} & 2009 & Rendimiento & Java EE & 3 capas & Experimento & Análisis automatizado; SGROAS incluye contraste no paramétrico \\ \midrule
\cite{inozemtseva2014} & 2014 & Cobertura & Varios & Varios & Experimento & Cobertura vs efectividad; SGROAS combina con integración \\ \midrule
\cite{brooke1996} & 1996 & Usabilidad & Varios & Varios & SUS & Instrumento estándar; SGROAS aplica con escala adjetiva \\ \midrule
\cite{pautasso2008} & 2008 & Web services & REST vs SOAP & REST & Comparación & REST vs big web services; SGROAS adopta REST \\ \midrule
\cite{fielding2002} & 2002 & Web & HTTP & REST & Diseño & Fundamentos REST; SGROAS sigue principios REST \\
\bottomrule
\end{longtable}

\section{Brecha identificada}
\label{sec:brecha}

Los trabajos revisados abordan dominios relacionados con transporte, arquitecturas web y evaluación empírica, pero presentan las siguientes limitaciones que SGROAS busca contribuir a cubrir:

\begin{enumerate}
  \item \textbf{Falta de sistemas web completos para cooperativas de transporte ecuatoriano:} la mayoría de trabajos se enfocan en microservicios \cite{alshuqayran2016,soldani2018,taibi2018,waseem2021} o en dominios distintos al transporte público en contextos latinoamericanos. SGROAS aporta un sistema completo (backend + frontend + BD) documentado con SRS ISO/IEC/IEEE 29148:2018.

  \item \textbf{Evaluación empírica incompleta:} varios trabajos reportan métricas de rendimiento sin intervalo de confianza \cite{jiang2009} o sin contraste estadístico formal. SGROAS incluye media, DT, IC 95\%, test inferencial (Wilcoxon) y tamaño de efecto (Cliff delta), siguiendo la guía de experimentación de Wohlin et al. \cite{wohlin2012}.

  \item \textbf{Ausencia de estrategia híbrida de acceso a datos documentada:} los trabajos de arquitectura web \cite{pautasso2008,fielding2002} describen patrones REST pero no documentan la separación CRUD-ORM / stored procedures con evidencia de seguridad (análisis estático de SQL dinámico). SGROAS documenta esta decisión formalmente (ADR-006) y la audita con scripts reproducibles.

  \item \textbf{Falta de reproducibilidad completa:} la mayoría de trabajos no publican datos crudos ni scripts de análisis. SGROAS cumple los principios FAIR \cite{wilkinson2016} con datos versionados, scripts reproducibles y DOI de Zenodo.
\end{enumerate}
